%==============================================================================
% Chapitre 9 - Humidité
%==============================================================================

\section{Introduction}

L'humidité atmosphérique correspond à la quantité de vapeur d'eau présente
dans l'air.

Son influence sur la balistique est réelle mais souvent surestimée.

Comparée à la température, à la pression atmosphérique ou au vent, son impact
reste généralement modéré.

Néanmoins, dans les applications exigeant une grande précision, elle peut
contribuer à modifier légèrement les trajectoires.

\begin{aretenir}

L'humidité influence la balistique principalement par son effet sur la
densité de l'air.

Son impact est généralement plus faible que celui de la température ou de la
pression.

\end{aretenir}

\section{La vapeur d'eau dans l'atmosphère}

L'air atmosphérique n'est pas constitué uniquement d'oxygène et d'azote.

Il contient également :

\begin{itemize}
\item du dioxyde de carbone ;
\item de l'argon ;
\item divers gaz traces ;
\item de la vapeur d'eau.
\end{itemize}

La quantité de vapeur d'eau varie continuellement selon les conditions
météorologiques.

\section{Humidité absolue et humidité relative}

Plusieurs grandeurs permettent de caractériser l'humidité.

La plus connue est l'humidité relative.

Elle représente le rapport entre :

\begin{itemize}
\item la quantité réelle de vapeur d'eau présente dans l'air ;
\item la quantité maximale que cet air pourrait contenir à la même
température.
\end{itemize}

Elle s'exprime généralement en pourcentage.

\begin{exemple}

Une humidité relative de 50\% signifie que l'air contient environ la moitié
de la vapeur d'eau qu'il pourrait théoriquement contenir à cette température.

\end{exemple}

\section{Pourquoi l'humidité influence-t-elle la densité ?}

Une idée reçue fréquente consiste à penser qu'un air humide est plus lourd
qu'un air sec.

En réalité, c'est généralement l'inverse.

La masse molaire de la vapeur d'eau est inférieure à celle de l'air sec.

Lorsqu'une partie des molécules de l'air est remplacée par de la vapeur
d'eau :

\begin{itemize}
\item la masse volumique du mélange diminue légèrement ;
\item la densité de l'air diminue.
\end{itemize}

\begin{aretenir}

À pression et température identiques, un air humide est généralement un peu
moins dense qu'un air sec.

\end{aretenir}

\section{Influence sur les trajectoires}

Une diminution de la densité de l'air entraîne généralement :

\begin{itemize}
\item une réduction de la traînée ;
\item une meilleure conservation de la vitesse ;
\item une légère augmentation de la portée.
\end{itemize}

L'effet reste toutefois limité dans la plupart des situations.

\begin{exemple}

Le passage d'une humidité relative de 20\% à 90\% produit généralement
des effets bien plus faibles qu'une variation importante de température ou de
pression atmosphérique.

\end{exemple}

\section{Humidité et vitesse du son}

La présence de vapeur d'eau modifie légèrement les propriétés acoustiques de
l'air.

La vitesse du son varie donc légèrement avec l'humidité.

Toutefois, dans la majorité des applications balistiques, cet effet demeure
secondaire devant l'influence de la température.

\section{Humidité et météorologie}

L'humidité constitue un indicateur important des conditions
météorologiques.

Elle influence notamment :

\begin{itemize}
\item la formation des nuages ;
\item le brouillard ;
\item les précipitations ;
\item la visibilité.
\end{itemize}

Ces phénomènes peuvent avoir des conséquences indirectes sur les opérations
de tir et les campagnes d'essais.

\section{Mesure de l'humidité}

L'humidité est généralement mesurée à l'aide d'un hygromètre.

Les stations météorologiques modernes fournissent souvent :

\begin{itemize}
\item la température ;
\item l'humidité relative ;
\item le point de rosée.
\end{itemize}

Ces informations peuvent être utilisées par certains calculateurs
balistiques.

\section{Application aux essais}

Lors des essais balistiques, l'humidité peut être enregistrée avec les autres
paramètres météorologiques.

Son influence étant relativement faible, elle n'est pas toujours prise en
compte dans les modèles simplifiés.

En revanche, elle est généralement intégrée dans les modèles les plus
précis.

\begin{simulation}

Dans de nombreux moteurs balistiques, l'humidité intervient dans le calcul de
la densité de l'air.

Son influence est souvent faible mais contribue à améliorer la fidélité du
modèle.

\end{simulation}

\section{Comparaison des influences atmosphériques}

Du point de vue balistique, l'importance relative des paramètres
atmosphériques est souvent la suivante :

\begin{enumerate}
\item vent ;
\item température ;
\item pression atmosphérique ;
\item humidité.
\end{enumerate}

Cette hiérarchie dépend naturellement :

\begin{itemize}
\item de la distance de tir ;
\item de la précision recherchée ;
\item du type de projectile.
\end{itemize}

\section{À retenir}

\begin{aretenir}

\begin{itemize}
\item L'humidité correspond à la quantité de vapeur d'eau présente dans
l'air.
\item Un air humide est généralement moins dense qu'un air sec.
\item L'humidité influence la traînée de manière indirecte.
\item Son effet est généralement plus faible que celui de la température
ou de la pression.
\item Elle peut néanmoins être prise en compte dans les modèles
balistiques de précision.
\end{itemize}

\end{aretenir}

